\section{Introduction}
\label{sec:introduction}

Large language models (LLMs) have rapidly become a practical means of generating unit tests from source code~\cite{schafer2023adaptive, chen2022codet, yuan2024chattester, lemieux2023codamosa, chen2021codex}, and an extensive recent literature has emerged around the technique~\cite{zhang2025llmunitsurvey}. Given a method signature, a modern LLM produces compilable, plausibly named tests in seconds; given the source body, it produces tests that exercise visible control-flow paths. The persistent weakness, well documented across the field~\cite{schafer2023adaptive, jiang2024survey, konstantinou2024oracles}, is the \emph{oracle problem}: the generated tests assert what the implementation does, not what it should do. Konstantinou et al.~\cite{konstantinou2024oracles} show empirically that LLM-generated oracles tend to encode the program's actual behaviour rather than its intended behaviour; when the implementation is buggy the model dutifully encodes the bug as the expected behaviour, and when the model is given only a signature it has no oracle at all and falls back to weak assertions such as ``does not throw''.

Formal specifications are the classical answer to the oracle problem~\cite{meyer1992design, jml, leavens2006design}, but their adoption is limited by the cost of writing them~\cite{matichuk_cost_2015, hall1990seven, zimmermann2019lightweight}. The empirical question this article addresses is therefore practical rather than theoretical: \emph{does giving a specification to an LLM measurably improve the unit tests it writes?} The hypothesis is that a specification---even one that is heuristically derived and incomplete---carries oracle information that is absent from the signature and largely redundant with the source body, and that this information is dense enough to shift LLM-generated tests from coverage-driven artefacts toward genuine behavioural checks.

The test-generation question is the headline empirical question because it has a clean answer under mutation testing, but it is one of several downstream uses of the same inferred specifications. The first is stronger LLM test oracles, the focus of the present article. The second is \emph{compositional verification of calling code}: a verifier that needs the contract of a callee in order to reason about a caller currently has either to verify the whole transitive dependency graph or to require manual contracts at every library boundary; inferred specifications act as the callee summaries that the caller's verification condition needs, so a change to a library entry implies a check of the calls that depend on it rather than a re-analysis of the entire program. The third is \emph{agentic programming}: an AI coding agent asked to interact with a library that postdates its training cut-off has no internal model of that library's contracts, and the same JML annotations function as a machine-readable contract the agent can consume to ground its reasoning. Throughout this article the inference machinery is treated as serving all three uses; the empirical evaluation isolates the first one because it is the use whose effect size mutation testing can measure within the budget of a single study.

The hypothesis is tested by holding the model fixed and varying only the input given to it. Four conditions are compared on the same 312-method corpus drawn from Apache Commons Lang: \textbf{P1}, the method signature alone; \textbf{P2}, signature with explicit guidance to cover edge cases; \textbf{P3}, signature together with a specification; \textbf{P4}, the full method source code as an upper-bound baseline. Test quality is measured by test count, compilation rate, pass rate, and mutation score (PIT~\cite{coles2016pitest}). P2 is included precisely to address the alternative hypothesis that any explicit guidance, not specifications in particular, would suffice: it does not---P2 generates the most tests but attains a lower mutation score than~P3, showing that test count is a poor proxy for oracle strength.

The methodological challenge is that 312 specifications cannot be written by hand within the budget of a single study. They are therefore inferred automatically by JML-Inferrer, a static analysis system developed alongside this study and described in Section~\ref{sec:approach}. The Java Modeling Language~(JML)~\cite{jml,leavens2006design} provides the contract notation used throughout. The JML-Inferrer is the means by which the empirical question becomes answerable at scale.

We make the following contributions:

\begin{enumerate}
    \item An affirmative empirical answer to the title question, with effect size: specifications increase the LLM's test count by 272\% (95\% bootstrap CI~$[245, 299]$) and mutation score by 40.7~percentage points relative to signatures alone ($p < 0.001$, Cohen's $d = 2.41$, 95\% bootstrap CI~$[2.08, 2.79]$).

    \item A breakdown of the effect by method category and by mutation operator, identifying where specifications carry the most oracle information (control-flow and utility methods) and where signature-only generation is already near-optimal (simple accessors).

    \item JML-Inferrer, the static analysis system used to produce the specifications, organised by weakest-precondition and strongest-postcondition reasoning and validated against OpenJML's Extended Static Checker.

    \item A probe-and-backport development methodology that uses an LLM to draft candidate JML, OpenJML's Extended Static Checker to filter discharged from undischarged clauses, and an agentic coding assistant to extend the inference engine until the discharged examples are emitted automatically. The methodology is what produced the heuristic repertoire underlying the experimental results, and is offered both as an account of how the system was developed and as a recipe for users who need to extend it (Section~\ref{subsec:probe-backport}).

    \item A replication package containing the source code, datasets, inferred specifications, generated test suites, and analysis scripts.
\end{enumerate}

Apache Commons Lang is used as the evaluation subject because it provides a mature, well-documented codebase with diverse method implementations and is widely adopted in software-engineering research~\cite{pacheco2007randoop, fraser2011evosuite}. The findings reported here motivate a planned follow-up study on service-boundary code (REST clients, RPC stubs), where the oracle problem is most acute and the effect reported here is expected to compound. The two remaining use cases identified above---compositional verification of callers and agentic programming over post-cut-off libraries---are the subject of separate ongoing work on specification distribution and on the agent-consumable contract format, and are mentioned here only to indicate that the inference machinery is not test-generation-specific.
